\subsection{Wye connections}
\label{sec:na-wye}

Figure~\ref{fig:na-wye} shows the three connections with a wye winding. All
three have the phase shift of 30\,$^{\circ}$ of Section~\ref{sec:na-terms}.

\begin{figure}[!ht]
  \centering
  % generated by north-american-practice/fig/connection_figures.py -- do not edit
\begin{tikzpicture}[x=0.27mm, y=-0.27mm, line cap=round, line join=round]
  \draw[hv] (38.0,130.0) -- (90.0,40.0);
  \node[font=\footnotesize, text=atInfo, anchor=center, inner sep=0.5pt] at (51.0,77.5) {480};
  \draw[hv] (90.0,40.0) -- (142.0,130.0);
  \node[font=\footnotesize, text=atInfo, anchor=center, inner sep=0.5pt] at (129.0,77.5) {480};
  \draw[hv] (142.0,130.0) -- (38.0,130.0);
  \node[font=\footnotesize, text=atInfo, anchor=center, inner sep=0.5pt] at (90.0,145.0) {480};
  \fill[atInfo] (38.0,130.0) circle (3.0);
  \node[font=\footnotesize, text=atInfo, anchor=center, inner sep=0.5pt] at (26.8,136.5) {H1};
  \fill[atInfo] (90.0,40.0) circle (3.0);
  \node[font=\footnotesize, text=atInfo, anchor=center, inner sep=0.5pt] at (90.0,27.0) {H2};
  \fill[atInfo] (142.0,130.0) circle (3.0);
  \node[font=\footnotesize, text=atInfo, anchor=center, inner sep=0.5pt] at (153.2,136.5) {H3};
  \draw[lv] (318.0,100.0) -- (258.0,100.0);
  \draw[lv] (318.0,100.0) -- (348.0,48.0);
  \draw[lv] (318.0,100.0) -- (348.0,152.0);
  \node[font=\footnotesize, text=atAccentFg, anchor=center, inner sep=0.5pt] at (288.0,87.0) {139};
  \fill[atFg] (318.0,100.0) circle (3.0);
  \draw[gnd] (318.0,103.0) -- (318.0,113.0);
  \draw[gnd] (311.0,113.0) -- (325.0,113.0);
  \draw[gnd] (313.5,117.0) -- (322.5,117.0);
  \draw[gnd] (316.0,121.0) -- (320.0,121.0);
  \node[font=\footnotesize, text=atFg, anchor=east, inner sep=0.5pt] at (306.0,120.0) {X0};
  \fill[atAccentFg] (258.0,100.0) circle (3.0);
  \node[font=\footnotesize, text=atAccentFg, anchor=center, inner sep=0.5pt] at (245.0,100.0) {X1};
  \fill[atAccentFg] (348.0,48.0) circle (3.0);
  \node[font=\footnotesize, text=atAccentFg, anchor=center, inner sep=0.5pt] at (354.5,36.8) {X2};
  \fill[atAccentFg] (348.0,152.0) circle (3.0);
  \node[font=\footnotesize, text=atAccentFg, anchor=center, inner sep=0.5pt] at (354.5,163.2) {X3};
  \draw[arr] (172.0,104.0) -- (204.0,104.0);
  \node[font=\footnotesize, text=atFgMuted, anchor=center, inner sep=0.5pt] at (188.0,122.0) {$-30$\textdegree};
  \node[font=\footnotesize, text=atFgMuted, anchor=center, inner sep=0.5pt] at (90.0,196.0) {480\,V $\Delta$};
  \node[font=\footnotesize, text=atFgMuted, anchor=center, inner sep=0.5pt] at (318.0,196.0) {240Y/139\,V};
  \node[font=\small\bfseries, text=atFg, anchor=west, inner sep=0.5pt] at (-6.0,18.0) {(a)};
  \draw[hv] (90.0,315.0) -- (38.0,345.0);
  \draw[hv] (90.0,315.0) -- (90.0,255.0);
  \draw[hv] (90.0,315.0) -- (142.0,345.0);
  \node[font=\footnotesize, text=atInfo, anchor=center, inner sep=0.5pt] at (70.5,341.3) {277};
  \fill[atFg] (90.0,315.0) circle (3.0);
  \node[font=\footnotesize, text=atDanger, anchor=center, inner sep=0.5pt] at (90.0,363.0) {H0 not connected};
  \fill[atInfo] (38.0,345.0) circle (3.0);
  \node[font=\footnotesize, text=atInfo, anchor=center, inner sep=0.5pt] at (26.8,351.5) {H1};
  \fill[atInfo] (90.0,255.0) circle (3.0);
  \node[font=\footnotesize, text=atInfo, anchor=center, inner sep=0.5pt] at (90.0,242.0) {H2};
  \fill[atInfo] (142.0,345.0) circle (3.0);
  \node[font=\footnotesize, text=atInfo, anchor=center, inner sep=0.5pt] at (153.2,351.5) {H3};
  \draw[lv] (258.0,315.0) -- (348.0,263.0);
  \node[font=\footnotesize, text=atAccentFg, anchor=center, inner sep=0.5pt] at (295.5,276.0) {240};
  \draw[lv] (348.0,263.0) -- (348.0,367.0);
  \node[font=\footnotesize, text=atAccentFg, anchor=center, inner sep=0.5pt] at (363.0,315.0) {240};
  \draw[lv] (348.0,367.0) -- (258.0,315.0);
  \node[font=\footnotesize, text=atAccentFg, anchor=center, inner sep=0.5pt] at (295.5,354.0) {240};
  \fill[atAccentFg] (258.0,315.0) circle (3.0);
  \node[font=\footnotesize, text=atAccentFg, anchor=center, inner sep=0.5pt] at (245.0,315.0) {X1};
  \fill[atAccentFg] (348.0,263.0) circle (3.0);
  \node[font=\footnotesize, text=atAccentFg, anchor=center, inner sep=0.5pt] at (354.5,251.8) {X2};
  \fill[atAccentFg] (348.0,367.0) circle (3.0);
  \node[font=\footnotesize, text=atAccentFg, anchor=center, inner sep=0.5pt] at (354.5,378.2) {X3};
  \draw[arr] (172.0,319.0) -- (204.0,319.0);
  \node[font=\footnotesize, text=atFgMuted, anchor=center, inner sep=0.5pt] at (188.0,337.0) {$-30$\textdegree};
  \node[font=\footnotesize, text=atFgMuted, anchor=center, inner sep=0.5pt] at (90.0,411.0) {480Y/277\,V};
  \node[font=\footnotesize, text=atFgMuted, anchor=center, inner sep=0.5pt] at (318.0,411.0) {240\,V $\Delta$};
  \node[font=\small\bfseries, text=atFg, anchor=west, inner sep=0.5pt] at (-6.0,233.0) {(b)};
  \draw[hv] (38.0,560.0) -- (90.0,470.0);
  \node[font=\footnotesize, text=atInfo, anchor=center, inner sep=0.5pt] at (51.0,507.5) {480};
  \draw[hv] (90.0,470.0) -- (142.0,560.0);
  \node[font=\footnotesize, text=atInfo, anchor=center, inner sep=0.5pt] at (129.0,507.5) {480};
  \draw[hv] (142.0,560.0) -- (38.0,560.0);
  \node[font=\footnotesize, text=atInfo, anchor=center, inner sep=0.5pt] at (90.0,575.0) {480};
  \fill[atInfo] (38.0,560.0) circle (3.0);
  \node[font=\footnotesize, text=atInfo, anchor=center, inner sep=0.5pt] at (26.8,566.5) {H1};
  \fill[atInfo] (90.0,470.0) circle (3.0);
  \node[font=\footnotesize, text=atInfo, anchor=center, inner sep=0.5pt] at (90.0,457.0) {H2};
  \fill[atInfo] (142.0,560.0) circle (3.0);
  \node[font=\footnotesize, text=atInfo, anchor=center, inner sep=0.5pt] at (153.2,566.5) {H3};
  \draw[lv] (318.0,530.0) -- (258.0,530.0);
  \draw[lv] (318.0,530.0) -- (348.0,478.0);
  \draw[lv] (318.0,530.0) -- (348.0,582.0);
  \node[font=\footnotesize, text=atAccentFg, anchor=center, inner sep=0.5pt] at (288.0,517.0) {120};
  \fill[atFg] (318.0,530.0) circle (3.0);
  \draw[gnd] (318.0,533.0) -- (318.0,543.0);
  \draw[gnd] (311.0,543.0) -- (325.0,543.0);
  \draw[gnd] (313.5,547.0) -- (322.5,547.0);
  \draw[gnd] (316.0,551.0) -- (320.0,551.0);
  \node[font=\footnotesize, text=atFg, anchor=east, inner sep=0.5pt] at (306.0,550.0) {X0};
  \fill[atAccentFg] (258.0,530.0) circle (3.0);
  \node[font=\footnotesize, text=atAccentFg, anchor=center, inner sep=0.5pt] at (245.0,530.0) {X1};
  \fill[atAccentFg] (348.0,478.0) circle (3.0);
  \node[font=\footnotesize, text=atAccentFg, anchor=center, inner sep=0.5pt] at (354.5,466.8) {X2};
  \fill[atAccentFg] (348.0,582.0) circle (3.0);
  \node[font=\footnotesize, text=atAccentFg, anchor=center, inner sep=0.5pt] at (354.5,593.2) {X3};
  \draw[arr] (172.0,534.0) -- (204.0,534.0);
  \node[font=\footnotesize, text=atFgMuted, anchor=center, inner sep=0.5pt] at (188.0,552.0) {$-30$\textdegree};
  \node[font=\footnotesize, text=atFgMuted, anchor=center, inner sep=0.5pt] at (90.0,626.0) {480\,V $\Delta$};
  \node[font=\footnotesize, text=atFgMuted, anchor=center, inner sep=0.5pt] at (318.0,626.0) {208Y/120\,V};
  \node[font=\small\bfseries, text=atFg, anchor=west, inner sep=0.5pt] at (-6.0,448.0) {(c)};
\end{tikzpicture}

  \caption{Connections with a wye winding: (a) 480\,V delta -- 240Y/139\,V,
  (b) 480Y/277\,V -- 240\,V delta with the primary neutral not connected,
  (c) 480\,V delta -- 208Y/120\,V for comparison.}
  \label{fig:na-wye}
\end{figure}

\subsubsection{480\,V delta -- 240Y/139\,V}
The secondary is a solidly grounded wye with 240\,V between lines and,
from Equation~\eqref{eq:na-uln}, 138.6\,V line to neutral. This voltage is
not a nominal system voltage of ANSI~C84.1~[5]; the neutral serves the
grounding of the system, not the supply of line-to-neutral loads.

\subsubsection{480Y/277\,V -- 240\,V delta}
On a 480Y/277\,V supply the primary is connected in wye with the neutral H0
left unconnected. If H0 is connected to the grounded supply neutral, the
delta secondary provides a closed path for zero-sequence current, and the
transformer becomes a grounding source for the 480\,V system: it carries
the zero-sequence part of the unbalanced 480\,V load and supplies current
into ground faults on the 480\,V system~[4]. The secondary is
operated in any of the ways of Section~\ref{sec:na-delta}.

\subsubsection{480\,V delta -- 208Y/120\,V}
The 208Y/120\,V secondary supplies 120\,V on all three phases from a solidly
grounded neutral and is the usual secondary when three-phase and 120\,V
loads are supplied together. A motor nameplated 230\,V operates at
$208/230 = 0.904$ of its rated voltage on this system, at the limit of the
$\pm$10\,\% voltage variation of NEMA~MG~1~[6]. Equipment rated 230 or
240\,V only is therefore supplied from a 240\,V secondary.
\FloatBarrier
